%% LyX 2.4.4 created this file.  For more info, see https://www.lyx.org/.
%% Do not edit unless you really know what you are doing.
\documentclass[10pt,english]{beamer}
\usepackage{mathptmx}
\usepackage[T1]{fontenc}
\usepackage[latin9]{inputenc}
\usepackage{amssymb}
\usepackage{graphicx}
\PassOptionsToPackage{normalem}{ulem}
\usepackage{ulem}

\makeatletter
%%%%%%%%%%%%%%%%%%%%%%%%%%%%%% Textclass specific LaTeX commands.
% this default might be overridden by plain title style
\newcommand\makebeamertitle{\frame{\maketitle}}%
% (ERT) argument for the TOC
\AtBeginDocument{%
  \let\origtableofcontents=\tableofcontents
  \def\tableofcontents{\@ifnextchar[{\origtableofcontents}{\gobbletableofcontents}}
  \def\gobbletableofcontents#1{\origtableofcontents}
}

%%%%%%%%%%%%%%%%%%%%%%%%%%%%%% User specified LaTeX commands.
\usetheme{Warsaw}
% or ...

\usepackage{pgfpages,pgfnodes,pgfarrows,pgfheaps}

\setbeamercovered{transparent}
% or whatever (possibly just delete it)
\setbeamercovered{invisible} %makes overlays invisible


%gets rid of navigation symbols
\setbeamertemplate{navigation symbols}{}
\usepackage{hyperref}

%%%setshandout pages per page and adds frames
%\usepackage{pgfpages}
%\pgfpagesuselayout{4 on 1}[a4paper, landscape, border shrink=7mm]
%\pgfpageslogicalpageoptions{1}{border code=\pgfusepath{stroke}}
%\pgfpageslogicalpageoptions{2}{border code=\pgfusepath{stroke}}
%\pgfpageslogicalpageoptions{3}{border code=\pgfusepath{stroke}}
%\pgfpageslogicalpageoptions{4}{border code=\pgfusepath{stroke}}

\makeatother

\usepackage{babel}
\begin{document}
%sets caption options. Requires \usepackage{caption} in preamble.
%\captionsetup[figure]{labelformat=empty, labelsep=none}
%the default beamer dimensions are 12.8 x 9.6cm.
\title{Circle Packing}
\subtitle{A Visual Introduction}
\author{James Ashe}
\titlegraphic{\includegraphics[height=3cm]{mhu_seal}}
\institute{Mars Hill University}
\date{}

\makebeamertitle

\beamerdefaultoverlayspecification{<+->}
\begin{frame}{What is a Circle Packing?}

\begin{columns}


\column{6.75cm}

\includegraphics[width=6.75cm]{unca_pics/0_owl1}<1|handout:0>\includegraphics[width=6.75cm]{unca_pics/0_owl2}<2|handout:0>\includegraphics[width=6.75cm]{unca_pics/0_owl3}<3|handout:0>\includegraphics[width=6.75cm]{unca_pics/0_owl4}<4>

\column{5.0cm}
\begin{block}{A collection of circles with a pattern such that}

\begin{enumerate}
\item [1.]<2->The circles have one and only one circle, $C_{i}$, associated
with each vertex $v_{i}$ in the pattern.
\item [2.]<3->Any two neighbors in the circle collection are \emph{tangent.}
\item [3.]<4->$\{C_{i},C_{j},C_{k}\}$ form a positively oriented triple
whenever $\{v_{i},v_{j},v_{k}\}$ form a positively oriented face
in the pattern.
\end{enumerate}
\end{block}
\end{columns}
\end{frame}

\begin{frame}{\textbf{What is a Circle Packing?}}

\uncover<3-5>{\textbf{Complex + Circles $=$ Geometry}}
\begin{columns}


\column{6.75cm }

\includegraphics[width=6.75cm]{unca_pics/owl1}<1|handout:0>\includegraphics[width=6.75cm]{unca_pics/owl2}<2|handout:0>\includegraphics[width=6.75cm]{unca_pics/owl3}<3|handout:0>\includegraphics[width=6.75cm]{unca_pics/owl4}<4|handout:0>\includegraphics[width=6.75cm]{unca_pics/owl5}<5|handout:0>
\end{columns}
\begin{itemize}
\item <5>Border conditions introduce flexibility.
\end{itemize}
\end{frame}

\begin{frame}{\textbf{Where do circle packings live}}

\begin{columns}


\column{3.5cm}
\begin{itemize}
\item <2->The plane, $\mathbb{C}$.
\item <3->The disc, $\mathbb{D}$.
\item <7->The sphere, $\mathbb{P}$.
\end{itemize}

\column{6.75cm }

\includegraphics[width=6.75cm]{unca_pics/owl_live0}<1|handout:0>\includegraphics[width=6.75cm]{unca_pics/owl_live1}<2|handout:0>\includegraphics[width=6.75cm]{unca_pics/owl_live2}<3|handout:0>\includegraphics[width=6.75cm]{unca_pics/owl_live3}<4|handout:0>\includegraphics[width=6.75cm]{unca_pics/owl_live4}<5|handout:0>\includegraphics[width=6.75cm]{unca_pics/owl_live5}<6|handout:0>\includegraphics[width=6.75cm]{unca_pics/owl_live6}<7|handout:0>\includegraphics[width=6.75cm]{unca_pics/owl_live7}<8|handout:0>
\end{columns}
\end{frame}

\begin{frame}{Moving between Geometries}

Via Stereographic Projections of the Riemann sphere.
\begin{itemize}
\item A M�bius transformation
\end{itemize}
\begin{columns}


\column{6cm}

\includegraphics[width=6cm]{unca_pics/owl_live5}<1|handout:0>\includegraphics[width=6cm]{unca_pics/3_owl1}<2|handout:0>\includegraphics[width=6cm]{unca_pics/3_owl2}<3|handout:0>

\column{6cm}

\includegraphics[width=6cm]{unca_pics/3_projection1}<1-2>\includegraphics[width=6cm]{unca_pics/3_projection2}<3->
\end{columns}
\end{frame}

\begin{frame}{KAT}

\begin{columns}


\column{12cm}
\begin{block}<2->{Koebe-Andreev-Thurston Theorem (KAT)}
 \emph{There exists a CP on the sphere for any finite triangulation
of the sphere. Furthermore, this CP is unique up to M�bius transformations.} 
\end{block}
\end{columns}
\begin{itemize}
\item [] \begin{center}When does a circle packing on $\mathbb{P}$ exist?

\includegraphics[height=4cm]{unca_pics/4_complex_to_sphere}<1-2|handout:0>\includegraphics[height=4cm]{unca_pics/3_owl1}<3|handout:0>\includegraphics[height=4cm]{unca_pics/3_owl2}<4-|handout:0>\end{center}
\end{itemize}
\end{frame}

\begin{frame}{Preliminaries for KAT Proof}

Hierarchy of a circle packing
\begin{columns}


\column{7cm}

\includegraphics[height=7.25cm]{\string"unca_pics/5_hierarchy_CP 1\string".eps}<1|handout:0>\includegraphics[height=7.25cm]{\string"unca_pics/5_hierarchy_CP 2\string".eps}<2|handout:0>\includegraphics[height=7.25cm]{\string"unca_pics/5_hierarchy_CP 3\string".eps}<3|handout:0>\includegraphics[height=7.25cm]{\string"unca_pics/5_hierarchy_CP 4\string".eps}<4>
\end{columns}
\end{frame}

\begin{frame}{Preliminaries for KAT Proof}

\begin{columns}


\column{6cm}

\vfill
\begin{block}{Monotonicity\emph{Given a face of circles at $v_{i}$, $v_{j}$, and
$v_{k}$ in $\mathbb{C}$ or $\mathbb{D}$. }}

\begin{itemize}
\item []\emph{The angle at $v_{i}$ is }\emph{\uline{strictly and continuously
}}\uline{decreasing}\emph{ wrt to the radius at $v_{i}$.}
\item []\emph{The angles at $v_{j}$ and $v_{k}$ are }\emph{\uline{strictly
and continuously }}\uline{increasing} \emph{wrt to the radius at $v_{i}$.}
\end{itemize}
\end{block}
\vfill

\column{5.5cm}

\includegraphics[height=5cm]{unca_pics/monotonicity/mono1}<1|handout:0>\includegraphics[height=5cm]{unca_pics/monotonicity/mono2}<2|handout:0>\includegraphics[height=5cm]{unca_pics/monotonicity/mono3}<3|handout:0>\includegraphics[height=5cm]{unca_pics/monotonicity/mono4}<4|handout:0>\includegraphics[height=5cm]{unca_pics/monotonicity/mono5}<5|handout:0>\includegraphics[height=5cm]{unca_pics/monotonicity/monoH1}<6|handout:0>\includegraphics[height=5cm]{unca_pics/monotonicity/monoH2}<7|handout:0>\includegraphics[height=5cm]{unca_pics/monotonicity/monoH3}<8|handout:0>\includegraphics[height=5cm]{unca_pics/monotonicity/monoH4}<9|handout:0>\includegraphics[height=5cm]{unca_pics/monotonicity/monoH5}<10|handout:0>\includegraphics[height=5cm]{unca_pics/monotonicity/monoH6}<11|handout:0>\includegraphics[height=5cm]{unca_pics/monotonicity/monoH7}<12|handout:0>\includegraphics[height=5cm]{unca_pics/monotonicity/monoH8}<13|handout:0>\includegraphics[height=5cm]{unca_pics/monotonicity/monoH9}<14|handout:0>\includegraphics[height=5cm]{unca_pics/monotonicity/monoH11}<15|handout:0>\includegraphics[height=5cm]{unca_pics/monotonicity/monoH14}<16|handout:0>\includegraphics[height=5cm]{unca_pics/monotonicity/flower1}<17|handout:0>\includegraphics[height=5cm]{unca_pics/monotonicity/flower2}<18|handout:0>\includegraphics[height=5cm]{unca_pics/monotonicity/flower3}<19|handout:0>\includegraphics[height=5cm]{unca_pics/monotonicity/flower4}<20|handout:0>\includegraphics[height=5cm]{unca_pics/monotonicity/flower5}<21|handout:0>\includegraphics[height=5cm]{unca_pics/monotonicity/flower6}<22|handout:0>\includegraphics[height=5cm]{unca_pics/monotonicity/flower5}<23|handout:0>\includegraphics[height=5cm]{unca_pics/monotonicity/flower3}<24|handout:0>\includegraphics[height=5cm]{unca_pics/monotonicity/flower1}<25|handout:0>\includegraphics[height=5cm]{unca_pics/monotonicity/flower7}<26|handout:0>\includegraphics[height=5cm]{unca_pics/monotonicity/flower8}<27|handout:0>\includegraphics[height=5cm]{unca_pics/monotonicity/flower9}<28-|handout:0>
\end{columns}
\begin{itemize}
\item <29->\emph{Packing condition: }Angle sum must be $2\pi n$.
\end{itemize}
\end{frame}

\begin{frame}{KAT Proof}

\begin{itemize}
\item []<2->Let $r_{i}$ be the radius at $v_{i}$. Define a \textbf{radii
assignment function},\textbf{\emph{ $\mathfrak{R}$}}\emph{
\[
\mathfrak{R}(r_{i})=\begin{cases}
\mbox{set \ensuremath{r_{i}}so the angle sum at \ensuremath{v_{i}}is \ensuremath{2\pi}} & \mbox{if }v_{i}\in\mbox{interior}\\
\mbox{unchanged} & \mbox{if }v_{i}\in\mbox{boundary}
\end{cases}
\]
}\vspace{-.5cm}
\end{itemize}
\begin{columns}


\column{4cm}

\includegraphics[height=5cm]{unca_pics/monotonicity/flower6}<2-|handout:0>

\column{.5cm}

\uncover<3->{$\longrightarrow$}

\column{4cm}

\hspace{-1cm}\includegraphics[height=5cm]{unca_pics/monotonicity/flower2}<4|handout:0>
\end{columns}
\end{frame}

\begin{frame}{KAT Proof}

\begin{itemize}
\item []<1->Let $r_{i}$ be the radius at $v_{i}$. Define a \textbf{radii
assignment function},\textbf{\emph{ $\mathfrak{R}$}}\emph{
\[
\mathfrak{R}(r_{i})=\begin{cases}
\mbox{set \ensuremath{r_{i}}so the angle sum at \ensuremath{v_{i}}is \ensuremath{2\pi}} & \mbox{if }v_{i}\in\mbox{interior}\\
\mbox{unchanged} & \mbox{if }v_{i}\in\mbox{boundary}
\end{cases}
\]
}\vspace{-.75cm}
\end{itemize}
\begin{columns}


\column{5cm}
\begin{itemize}
\item Start in $\mathbb{C}$ or $\mathbb{D}$, so we have

\begin{itemize}
\item Border circles
\item Monotonicity
\end{itemize}
\end{itemize}
\begin{block}<2->{Lemma}
 Given two radii collection for a pattern, such that $R'\leq R$
then $\mathfrak{R}(R')\leq\mathfrak{R}(R)$.
\end{block}

\column{7cm}
\begin{itemize}
\item [] \uncover<3->{\includegraphics[width=3cm]{unca_pics/6_flower1}}\uncover<5->{\includegraphics[width=3cm]{unca_pics/6_flower2}}
\item [] \vspace{-.3cm} \uncover<4->{\includegraphics[width=3cm]{unca_pics/6_flower3}}\uncover<6->{\includegraphics[width=3cm]{unca_pics/6_flower4}}
\end{itemize}
\end{columns}
\end{frame}

\begin{frame}{KAT Proof}

\begin{itemize}
\item []<1->Let $r_{i}$ be the radius at $v_{i}$. Define a \textbf{radii
assignment function},\textbf{\emph{ $\mathfrak{R}$}}\emph{
\[
\mathfrak{R}(r_{i})=\begin{cases}
\mbox{set \ensuremath{r_{i}}so the angle sum at \ensuremath{v_{i}}is \ensuremath{2\pi}} & \mbox{if }v_{i}\in\mbox{interior}\\
\mbox{unchanged} & \mbox{if }v_{i}\in\mbox{boundary}
\end{cases}
\]
}\vspace{-.75cm}
\end{itemize}
\begin{columns}


\column{12cm}
\begin{block}<1->{Lemma}
 Given two radii collection for a pattern, such that $R'\leq R$
then $\mathfrak{R}(R')\leq\mathfrak{R}(R)$.
\end{block}
\begin{itemize}
\item []<2->So $\mathfrak{R}(R)\leq R\iff\mathfrak{R}^{k+1}(R')\leq\mathfrak{R}^{k}(R)$
\item []<3->and \textrm{$\mathfrak{R}(R)\leq R\iff$angle sums of $R\leq2\pi$.}
\item []<4->and \textrm{$\mathfrak{R}^{k+1}(R)\leq\mathfrak{R}^{k}(R)\iff$angle
sums of $\mathfrak{R}^{k}(R)\leq2\pi$.}
\end{itemize}
\end{columns}
\begin{itemize}
\item <5->Now choose some really big radii; so that $\mathfrak{R}(R_{\infty})<R_{\infty}$,
and it converges.
\item <6->$\dots$with some details. 
\end{itemize}
\end{frame}

\begin{frame}{KAT Proof}

\begin{columns}


\column{5.5cm}
\begin{block}{Theorem}
\emph{ There exists a unique CP for any simply connected pattern
with a collection of border circles in $\mathbb{C}$ or $\mathbb{D}$.}
\end{block}
\begin{itemize}
\item <2->Maximality argument$\Rightarrow$uniqueness.
\item <3->To get on the sphere$\dots$
\end{itemize}
\begin{block}<8->{Koebe-Andreev-Thurston Theorem (KAT)}
 \emph{There exists a CP on the sphere for any finite triangulation
of the sphere. Furthermore, this CP is unique up to M�bius transformations.} 
\end{block}

\column{5.5cm}

\includegraphics[width=5.5cm]{unca_pics/6_flower2}<2>\includegraphics[width=5.5cm]{unca_pics/6_flower1}<3>\includegraphics[width=5.5cm]{unca_pics/owl_live0}<4>\includegraphics[width=5.5cm]{unca_pics/owl_live2}<5>\includegraphics[width=5.5cm]{unca_pics/owl_live5}<6>\includegraphics[width=5.5cm]{unca_pics/owl_live7}<7->
\end{columns}
\end{frame}

\begin{frame}{Circle Packings as Discrete Analytical Function}

\begin{itemize}
\item <2->Approximation and computation of complex functions. 
\item <3->Discrete analogues of complex functions.
\end{itemize}
\only<4>{\includegraphics[width=10cm]{unca_pics/map1/map1}}\only<5>{\includegraphics[width=10cm]{unca_pics/map1/map2}}\only<6>{\includegraphics[width=10cm]{unca_pics/map1/map3}}\only<7>{\includegraphics[width=10cm]{unca_pics/map1/map4}}\only<8>{\includegraphics[width=10cm]{unca_pics/map1/map5}}\only<9>{\includegraphics[width=10cm]{unca_pics/map1/map6}}\uncover<10->{\includegraphics[width=10cm]{unca_pics/map1/map7}}
\begin{itemize}
\item []<11->Analytic functions, \emph{map infinitesimal circles to infinitesimal
circles.}
\item []<12->Circle packing, \emph{maps real circles to real circles}. 
\end{itemize}
\end{frame}

\begin{frame}{Riemann Mapping Theorem}

\only<1-1>{
\begin{block}{Classical RMT}
\emph{If $R$ is a non-empty simply connected proper subset of $\mathbb{C}$,
then there exists an essentially unique mapping from $R$ onto the
unit disc. }
\end{block}
}

\only<2->{
\begin{block}{Discrete RMT}
\emph{If $C$ is a simply connected finite circle packing in $\mathbb{C}$,
then there exists an essentially unique maximal packing of with the
same pattern in $\mathbb{D}$. }
\end{block}
}

\only<1>{\includegraphics[width=11cm]{unca_pics/RMT/1_RMT1a}}\only<2>{\includegraphics[width=11cm]{unca_pics/RMT/1_RMT2}}\only<3>{\includegraphics[width=11cm]{unca_pics/RMT/1_RMT3}}\only<4>{\includegraphics[width=11cm]{unca_pics/RMT/1_RMT4}}\only<5>{\includegraphics[width=11cm]{unca_pics/RMT/1_RMT5}}\only<6>{\includegraphics[width=11cm]{unca_pics/RMT/1_RMT6}}\only<7>{\includegraphics[width=11cm]{unca_pics/RMT/1_RMT8}}\only<8>{\includegraphics[width=11cm]{unca_pics/RMT/1_RMT10}}\only<9>{\includegraphics[width=11cm]{unca_pics/RMT/1_RMT12}}\only<10>{\includegraphics[width=11cm]{unca_pics/RMT/1_RMT14}}\only<11>{\includegraphics[width=11cm]{unca_pics/RMT/1_RMT15}}\only<12>{\includegraphics[width=11cm]{unca_pics/RMT/2_R1}}\only<13>{\includegraphics[width=11cm]{unca_pics/RMT/2_R2}}\only<14>{\includegraphics[width=11cm]{unca_pics/RMT/2_R3}}\only<15>{\includegraphics[width=11cm]{unca_pics/RMT/2_R6}}\only<16>{\includegraphics[width=11cm]{unca_pics/RMT/2_R9}}\only<17>{\includegraphics[width=11cm]{unca_pics/RMT/2_R12}}\only<18>{\includegraphics[width=11cm]{unca_pics/RMT/2_R14}}\only<19>{\includegraphics[width=11cm]{unca_pics/RMT/3R1}}\only<20>{\includegraphics[width=11cm]{unca_pics/RMT/3R2}}\only<21>{\includegraphics[width=11cm]{unca_pics/RMT/3R3}}\only<22>{\includegraphics[width=11cm]{unca_pics/RMT/3R4}}\only<23>{\includegraphics[width=11cm]{unca_pics/RMT/3R5}}\only<24>{\includegraphics[width=11cm]{unca_pics/RMT/3R7}}\only<25>{\includegraphics[width=11cm]{unca_pics/RMT/3R9}}\only<26>{\includegraphics[width=11cm]{unca_pics/RMT/3R12}}\only<27>{\includegraphics[width=11cm]{unca_pics/RMT/1_RMT1a}}
\begin{itemize}
\item <26->More circles, and the maps converge to the classical mapping.
\end{itemize}
\end{frame}

\begin{frame}{Discrete \underline{\hspace{3cm}} Function}

\begin{itemize}
\item Discrete Schwarz Lemma
\item Discrete Blaschke products
\item Discrete polynomials
\item Discrete rational functions, etc.
\end{itemize}
\begin{block}<2->{Theorem}
\emph{Discrete analytic functions converge under refinement to their
classical counterparts. }
\end{block}
\begin{itemize}
\item []<3->\textcolor{blue}{{[}He and Schramm{]}} The KAT theorem implies
the classical Riemann Mapping Theorem 
\end{itemize}
\end{frame}

\begin{frame}{Branching}

\begin{columns}


\column{5cm}
\begin{exampleblock}<1->{Flowers with angle sums of}

\begin{itemize}
\item <2-8>$2\pi$, a (trivially) branched flower.
\item <9-15>$4\pi$, a branched flower.
\item <16->$2.75\pi$, an ``incoherent'' flower
\end{itemize}
\end{exampleblock}

\column{6cm}

\includegraphics[width=5cm]{../Thesis_defense/aux_pics/cp_5587/cpShot-046}<2|handout:0>\includegraphics[width=5cm]{../Thesis_defense/aux_pics/cp_5587/cpShot-047}<3|handout:0>\includegraphics[width=5cm]{../Thesis_defense/aux_pics/cp_5587/cpShot-048}<4|handout:0>\includegraphics[width=5cm]{../Thesis_defense/aux_pics/cp_5587/cpShot-049}<5|handout:0>\includegraphics[width=5cm]{../Thesis_defense/aux_pics/cp_5587/cpShot-050}<6|handout:0>\includegraphics[width=5cm]{../Thesis_defense/aux_pics/cp_5587/cpShot-051}<7|handout:0>\includegraphics[width=5cm]{../Thesis_defense/aux_pics/cp_5587/cpShot-052}<8|handout:0>\includegraphics[width=5cm]{../Thesis_defense/aux_pics/cp_5587/cpShot4pi-053}<9|handout:0>\includegraphics[width=5cm]{../Thesis_defense/aux_pics/cp_5587/cpShot4pi-054}<10|handout:0>\includegraphics[width=5cm]{../Thesis_defense/aux_pics/cp_5587/cpShot4pi-055}<11|handout:0>\includegraphics[width=5cm]{../Thesis_defense/aux_pics/cp_5587/cpShot4pi-056}<12|handout:0>\includegraphics[width=5cm]{../Thesis_defense/aux_pics/cp_5587/cpShot4pi-057}<13|handout:0>\includegraphics[width=5cm]{../Thesis_defense/aux_pics/cp_5587/cpShot4pi-058}<14|handout:0>\includegraphics[width=5cm]{../Thesis_defense/aux_pics/cp_5587/cpShot4pi-059}<15|handout:0>\includegraphics[width=5cm]{../Thesis_defense/aux_pics/cp_5587/cpShot2-75pi-067}<16|handout:0>\includegraphics[width=5cm]{../Thesis_defense/aux_pics/cp_5587/cpShot2-75pi-068}<17|handout:0>\includegraphics[width=5cm]{../Thesis_defense/aux_pics/cp_5587/cpShot2-75pi-069}<18|handout:0>\includegraphics[width=5cm]{../Thesis_defense/aux_pics/cp_5587/cpShot2-75pi-070}<19|handout:0>\includegraphics[width=5cm]{../Thesis_defense/aux_pics/cp_5587/cpShot2-75pi-071}<20|handout:0>\includegraphics[width=5cm]{../Thesis_defense/aux_pics/cp_5587/cpShot2-75pi-072}<21|handout:0>\includegraphics[width=5cm]{../Thesis_defense/aux_pics/cp_5587/cpShot2-75pi-073}<22|handout:0>\includegraphics[width=5cm]{../Thesis_defense/aux_pics/cp_5587/cpShot2-75pi-074}<23|handout:0>\uncover<24->{\includegraphics[width=5cm]{../Thesis_defense/aux_pics/cp_5587/cpShot2-75pi-075}}
\end{columns}
\end{frame}

\begin{frame}{Branching}

\begin{columns}


\column{5cm}
\begin{exampleblock}<1->{Interior flowers with angle sums of}

\begin{itemize}
\item <2-21>$4\pi$, a branched flower.
\item <22->$3.66\pi$, an ``incoherent'' flower
\end{itemize}
\end{exampleblock}
\begin{center} \uncover<12-23>{\includegraphics[height=4cm]{unca_pics/branched_sph}
\begin{itemize}
\item Given some combinatorial restrictions!}
\end{itemize}
\end{center}

\column{6cm}

\includegraphics[width=6cm]{unca_pics/holonomy2/1branch1}<1|handout:0>\includegraphics[width=6cm]{unca_pics/holonomy2/1branch2}<2>\includegraphics[width=6cm]{unca_pics/holonomy2/1branch3}<3>\includegraphics[width=6cm]{unca_pics/holonomy2/1branch4}<4>\includegraphics[width=6cm]{unca_pics/holonomy2/1branch5}<5>\includegraphics[width=6cm]{unca_pics/holonomy2/1branch6}<6>\includegraphics[width=6cm]{unca_pics/holonomy2/1branch7}<7>\includegraphics[width=6cm]{unca_pics/holonomy2/1branch8}<8>\includegraphics[width=6cm]{unca_pics/holonomy2/1branch9}<9>\includegraphics[width=6cm]{unca_pics/holonomy2/1branch10}<10>\includegraphics[width=6cm]{unca_pics/holonomy2/1branch11}<11>\includegraphics[width=6cm]{unca_pics/holonomy2/1branch12}<12>\includegraphics[width=6cm]{unca_pics/holonomy2/1branch13}<13>\includegraphics[width=6cm]{unca_pics/holonomy2/1small1}<14>\includegraphics[width=6cm]{unca_pics/holonomy2/1small2}<15>\includegraphics[width=6cm]{unca_pics/holonomy2/1small3}<16>\includegraphics[width=6cm]{unca_pics/holonomy2/1small4}<17>\includegraphics[width=6cm]{unca_pics/holonomy2/1small5}<18>\includegraphics[width=6cm]{unca_pics/holonomy2/1small6}<19>\includegraphics[width=6cm]{unca_pics/holonomy2/1small7}<20>\includegraphics[width=6cm]{unca_pics/holonomy2/1small8}<21>\includegraphics[width=6cm]{unca_pics/holonomy2/2small1}<22>\includegraphics[width=6cm]{unca_pics/holonomy2/1small2}<23>\includegraphics[width=6cm]{unca_pics/holonomy2/2small3}<24>\includegraphics[width=6cm]{unca_pics/holonomy2/2small4}<25>\includegraphics[width=6cm]{unca_pics/holonomy2/2small5}<26>\includegraphics[width=6cm]{unca_pics/holonomy2/2small6}<27>\includegraphics[width=6cm]{unca_pics/holonomy2/2small7}<28>\includegraphics[width=6cm]{unca_pics/holonomy2/2small8}<29>\includegraphics[width=6cm]{unca_pics/holonomy2/2small9}<30>\includegraphics[width=6cm]{unca_pics/holonomy2/2small10}<31>\includegraphics[width=6cm]{unca_pics/holonomy2/2branch1}<32>\includegraphics[width=6cm]{unca_pics/holonomy2/2branch2}<33>\includegraphics[width=6cm]{unca_pics/holonomy2/2branch3}<34>\includegraphics[width=6cm]{unca_pics/holonomy2/2branch4}<35>\includegraphics[width=6cm]{unca_pics/holonomy2/2branch5}<36>\includegraphics[width=6cm]{unca_pics/holonomy2/2branch6}<37>\includegraphics[width=6cm]{unca_pics/holonomy2/2branch7}<38>\includegraphics[width=6cm]{unca_pics/holonomy2/2branch8}<39>\uncover<40->{\includegraphics[width=6cm]{unca_pics/holonomy2/2branch9}}
\end{columns}
\end{frame}

\begin{frame}{Fractional Branching}

\begin{columns}


\column{6cm}

Choose an interior face $f=\{v_{1},v_{2},v_{3}\}$. Is there a way
to split the extra angle amongst the vertices to get a branched packing? 
\begin{block}<4->{Theorem}

$\iff$the extra angle at $v_{i}$ is $2(\angle v_{i}v_{j}$$v_{k})$
for $v_{i},v_{j},v_{k}\in f$.
\end{block}

\begin{block}<11->{Theorem}

If the pattern is simply connected then there (usually) exists a fractional
branching on very interior face\textbf{.}
\end{block}
\begin{itemize}
\item <12->Proof using Brouwer fixed point theorem. 
\end{itemize}

\column{5cm}

\includegraphics[width=5cm]{unca_pics/fractional_data2}<2>\includegraphics[width=5cm]{unca_pics/fractional_data1}<3-4>\includegraphics[width=5cm]{unca_pics/rotating_trinagles1}<5>\includegraphics[width=5cm]{unca_pics/rotating_trinagles2}<6>\includegraphics[width=5cm]{unca_pics/rotating_trinagles3}<7>\includegraphics[width=5cm]{../Thesis_defense/pictures/FB_simple_unfrac}<8>\includegraphics[width=5cm]{../Thesis_defense/pictures/FB_simple_frac}<9>\includegraphics[width=5cm]{../Thesis_defense/pictures/FB_simple_frac_detail}<10>\includegraphics[width=5cm]{../Thesis_defense/pictures/FB_simple_frac_detail_b}<11->
\end{columns}
\end{frame}

\begin{frame}{Thomson Problem}

\begin{columns}


\column{6cm}

How to find the min. energy distribution of $N$ electrons on the
surface of a sphere?
\[
E=\sum_{i<j}\frac{1}{\left|p_{i}-p_{j}\right|}
\]

\begin{itemize}
\item <2-> A CP exists for any triangulation.
\item <3-> The CP can be computed by \href{http://www.math.utk.edu/~kens/CirclePack/}{\textcolor{blue}{\tt{CirclePack}}}.
\item <4-> As $N\rightarrow\infty$, a CP analogue will converge towards
its continuous counterpart.
\end{itemize}
\begin{exampleblock}<4->{}

\begin{itemize}
\item <5->[1.] Triangulate existing data pts.
\item <6->[2.] Compute using \tt{CirclePack}.
\item <7->[3.] Normalize the packing.
\end{itemize}
\end{exampleblock}

\column{5cm}

\includegraphics[width=5cm]{unca_pics/thomson_pics/Thomson_410_mesh_thermal}<1>\includegraphics[width=5cm]{unca_pics/thomson_pics/Thomson_410_3d_pts}<2-5>\includegraphics[width=5cm]{unca_pics/thomson_pics/cpShot-001}<6-7>\includegraphics[width=5cm]{unca_pics/thomson_pics/data_and_chart_for_printing3}<8>\includegraphics[width=5cm]{unca_pics/thomson_pics/ladder4}<9>\includegraphics[width=5cm]{unca_pics/thomson_pics/ladder3}<10>\includegraphics[width=5cm]{unca_pics/thomson_pics/ladder2}<11>\includegraphics[width=5cm]{unca_pics/thomson_pics/ladder1}<12>
\end{columns}
\end{frame}

\begin{frame}{}

All images created with ${\tt CirclePack}$

\href{http://www.math.utk.edu/~kens/CirclePack/}{\textcolor{blue}{http://www.math.utk.edu/$\sim$kens/CirclePack/}}

\begin{center}
\includegraphics[height=6.5cm]{../AMS_sectional_Knoxville_2014/pics/AMS_Knoxville_2014_pics/CP_screenshot_3_09}
\par\end{center}

\end{frame}

\begin{frame}[allowframebreaks]{References}

\nocite{*}\bibliographystyle{plain}
\bibliography{uncaa_jr}

\end{frame}

\begin{frame}{}

\begin{center}
\Large\textbf{Thanks!}
\par\end{center}
\end{frame}

\end{document}
